%% =====================================================================
%%  B4-T-14-01 -- what B4 contributes to "The Automatic Response"
%%  -------------------------------------------------------------------
%%  LAYER A: APPEND-ONLY. Written once, then left alone. Material found
%%  only here sits in the `unique` slot and is protected by rule R10.
%% =====================================================================
%% @kind: source
%% @id: B4-T-14-01
%% @book: B4
%% @topic: T-14-01
%% @locator: intro, pp. 1--12; ch. 1, pp. 13--42; epilogue
%% @status: distilled
%% @unique: yes
%% @integrated: yes
%% @updated: 2026-09-19

\dossier{B4}{T-14-01}{\bookshort{B4}}{intro, pp. 1--12; ch. 1, pp. 13--42; epilogue}

\begin{distilled}
  \item Compliance is frequently produced by a single trigger feature rather than by an assessment of the request; the response pattern is fixed and the cue that releases it is separable from the substance.
  \item The shortcut is retained because full evaluation is too costly to perform on every ordinary request, so it is correct often enough to be worth keeping.
  \item Each subsequent principle in the book is one such feature --- obligation, consistency, crowd, warmth, credential, rarity --- exploited by producing the feature rather than the condition.
  \item The accelerating case: cues now arrive faster than they can be evaluated, so more decisions are made on the shortcut than before.
\end{distilled}

\begin{unique}
  \item The trigger-feature model itself: the claim that a whole compliance routine can be released by one manufactured cue, which is the causal layer under the tactics in B1 and the laws in B3.
  \item The economic justification for keeping the shortcut, and the resulting asymmetry of the defence.
\end{unique}
